% VI36-DETAILED-PROBLEM-16
\begin{problem}[VI.36.P16: Homogenization and dehomogenization]\label{prob:vi36-16}
Let \(f\in k[t_1,\ldots,t_n]\) have total degree at most \(d\). Define its degree-\(d\) homogenization
\[
F(x_0,\ldots,x_n)=x_0^d f(x_1/x_0,\ldots,x_n/x_0).
\]
Prove that \(F\) is homogeneous of degree \(d\), and that dehomogenizing at \(x_0=1\) recovers \(f\).

\textbf{Provenance.} Canonical VI/36 extension.
\end{problem}

\ifdefined\FullProblemDossiers
\begin{solution}
Write
\[
f=\sum_\alpha c_\alpha t^\alpha,
\qquad
|\alpha|\le d.
\]
Then
\[
F
=
\sum_\alpha
c_\alpha
x_0^{d-|\alpha|}
x_1^{\alpha_1}\cdots x_n^{\alpha_n}.
\]
Each monomial has total degree
\[
(d-|\alpha|)+|\alpha|=d.
\]
Hence
\[
\boxed{F\text{ is homogeneous of degree }d}.
\]

Setting \(x_0=1\) gives
\[
F(1,x_1,\ldots,x_n)
=
\sum_\alpha
c_\alpha x_1^{\alpha_1}\cdots x_n^{\alpha_n}
=
f(x_1,\ldots,x_n).
\]
Thus dehomogenization on \(U_0\) recovers the original affine polynomial exactly.

If \(d=\deg f\), this is the usual homogenization used to study the projective completion of the affine zero locus.
The full scheme-theoretic closure statement is reserved for VI/37.
\end{solution}
\fi
